\section{Parametrización del proceso inverso}

\subsection{Modelado de la distribución de transferencia inversa}

Dado que el posterior real $q(x_{t-1}|x_t, x_0)$ es una distribución gaussiana, naturalmente modelamos también la transferencia inversa $p_\theta(x_{t-1}|x_t)$ como una distribución gaussiana:

$$p_\theta(x_{t-1}|x_t) = \mathcal{N}\big(x_{t-1};\, \mu_\theta(x_t, t),\; \sigma_t^2 I\big)$$

Donde:
\begin{itemize}
    \item $\mu_\theta(x_t, t)$: la media predicha por una red neuronal, tomando como entrada $x_t$ y el paso de tiempo $t$
    \item $\sigma_t^2$: la varianza, fijada en $\tilde{\beta}_t$ (igual que en el posterior real), \textbf{sin necesidad de aprenderla}
\end{itemize}

\begin{knowledgebox}{¿Por qué no es necesario aprender la varianza?}
Dado que queremos minimizar $D_{\text{KL}}(q(x_{t-1}|x_t,x_0) \| p_\theta(x_{t-1}|x_t))$, y la divergencia de KL entre dos distribuciones gaussianas se simplifica a la norma L2 de la diferencia de medias cuando las varianzas son iguales. Por lo tanto, fijar la varianza de $p_\theta$ en $\tilde{\beta}_t$, igual que en $q$, hace que el objetivo de entrenamiento se reduzca únicamente a ajustar la media. El artículo Improved DDPM explora esquemas para aprender la varianza, pero en DDPM básico es suficiente con fijarla.
\end{knowledgebox}

\subsection{Objetivo de entrenamiento simplificado}

Sustituyendo la forma gaussiana en la divergencia de KL, el término de coincidencia de des ruido $L_{t-1}$ se simplifica a:

$$L_{t-1} = \frac{1}{2\sigma_t^2} \left\| \tilde{\mu}_t(x_t, x_0) - \mu_\theta(x_t, t) \right\|^2 + C$$

Donde $C$ es una constante independiente de $\theta$. Por lo tanto, el objetivo de entrenamiento consiste en que la media predicha por la red neuronal $\mu_\theta(x_t, t)$ se acerque lo más posible a la media del posterior real $\tilde{\mu}_t(x_t, x_0)$.

\begin{importantbox}{Intuición del entrenamiento}
El proceso de entrenamiento puede describirse brevemente como: (1) muestrear un dato limpio $x_0$; (2) muestrear el paso de tiempo $t$; (3) muestrear datos con ruido $x_t$ mediante la distribución de salto hacia adelante; (4) predecir la media $\mu_\theta(x_t, t)$ con una red neuronal; (5) calcular la pérdida L2 entre la media predicha y la real, y realizar el retroceso.
\end{importantbox}

En el entrenamiento práctico, el coeficiente $\frac{1}{2\sigma_t^2}$ equivale a un término de peso dependiente del tiempo $t$. Ho et al. (2020) descubrieron que, en la práctica, \textbf{omitir este peso} (es decir, usar pesos uniformes para todos los pasos de tiempo) funciona muy bien e incluso mejor. Este es un hallazgo empírico importante del artículo sobre DDPM.

\subsection{Resumen de la sección}

Al modelar la transferencia inversa como una distribución gaussiana y fijar la varianza, el objetivo de entrenamiento se simplifica a un problema de regresión L2 de medias. El término de peso suele omitirse en la práctica para simplificar el entrenamiento.

%% ========== Sección 5: Tres objetivos de predicción ==========


